\begin{claim}{C.64}{Proved}{Compact local control and equivariant disk localization at a fixed frontier}
Let $f_t$ and $G$ be the original based ambient-null plane loop and
its generated group, and let $H$ commute with every $f_t$. Let $W$
be an $H$-invariant union of open $G$-orbits. There is a compact
path-connected labelled family $\mathcal K$ of actual global
$H$-centralizing homeomorphisms, containing the identity, such that
every member preserves each original open orbit and fixes the
$G$-fixed set, and
\[
 x\in\operatorname{int}_{O}
       \bigl(\mathcal K\mathcal K^{-1}(x)\bigr)
       \qquad(x\in O\subset W).                         \tag{C.64.1}
\]
All original source-cover labels are retained in this family.

Now suppose that $\partial W$ is $G$-fixed and that the full local
centralizer on every $O\subset W$ has dimension zero. One can
choose arbitrarily small closed embedded disks $D_i\subset O_i$,
with points $x_i\in D_i$, for which
\[
 D_i\subset\mathcal K\mathcal K^{-1}(x_i),               \tag{C.64.2}
\]
and every effective stabilizer translate of $D_i$ is disjoint
from it. Take any countable such family, with the $D_i$
representing pairwise distinct $H$-orbits of disks, whose
distinct $H$-translates are pairwise disjoint and locally
finite in $W$.
There is an actual continuous group embedding, with continuous
inverse on its image,
\[
 \mathfrak L:\prod_i\operatorname{Homeo}(D_i\mathbin{\mathrm{rel}}\partial D_i)
 \longrightarrow
 C_{\operatorname{Homeo}(\R^2)}(H),                       \tag{C.64.3}
\]
obtained by conjugating each disk map over its $H$-translates
and using the identity elsewhere. Its image is a contractible
subgroup. Both the maps and inverses are jointly continuous,
including at every accumulation on $\partial W$, uniformly
over this whole product of disk maps. In particular every
based loop in this subgroup has an explicitly prescribed
based contraction in the same global centralizer.
\end{claim}

\begin{proof}
\emph{The exact path group.}
Use every original universal cover $p_O:X_O\to O$ of the selected
open orbits, with the original based lifted loop. The open orbits
are countable, by a countable basis of the plane. The closed
compatible tuple construction (C.63.10) therefore gives a Polish
group of global homeomorphisms and their labelled original lifts.
It imposes $p_O\widetilde k^O=kp_O$, the inverse equations,
global commutation with $H$, and commutation with the retained
lifts of the full local centralizers. This construction uses
no assumption on their dimension. The original loop belongs
to this group: its commutators with fixed operator lifts are
deck-valued and initially trivial, and its based lifts close
by C.32. Let $L$ be the closure of the subgroup generated by
that labelled loop. Its global values preserve every selected
open orbit and fix every $G$-fixed point, since these are closed
conditions for maps and inverses.

Write $\mathscr P(L)$ for the group of continuous paths
$\sigma:[0,1]\to L$ with $\sigma(0)=1$, with pointwise
multiplication and the compact-open path topology. A complete
compatible metric on $L$ gives a complete uniform path metric.
Uniform limits are continuous and based. The space is separable:
uniform-continuity scales and finitely many values in a countable
dense set give a countable neighborhood coding; choose one actual
path from each nonempty code. Multiplication and inversion are
continuous uniformly on compact parameter sets. Thus
$\mathscr P(L)$ is a Polish group.

For $x\in O$, the continuous action by endpoint evaluation is
transitive. A finite original word carrying $x$ to a given point
of $O$ is the endpoint of a based word path: replace every
original time $t_j$ by $s t_j$, preserving the order and inverse
sign of every letter. These are paths in $L$, with their original
source lifts. The precise orbit-open-map theorem used in C.30,
\cite[Theorem~3.5]{C_GomezPinedoUzcategui2017}, consequently makes
\[
 e_x:\mathscr P(L)\longrightarrow O,\qquad
                 e_x(\sigma)=\sigma(1)x                 \tag{C.64.4}
\]
an open surjection. Its source is Polish and its target is the
nonmeagre Hausdorff surface $O$, as required there.

\emph{Compact lifting of a compact target.}
We give the needed elementary consequence of this open map.
Let $p:X\to Y$ be a continuous open surjection from a completely
metrizable space to a metric space. For an open $V\subset X$
and a compact $D\subset p(V)$, there is a compact
$S\subset V$ with $p(S)=D$.

Here is the full construction, including the equality of images.
For each $y\in D$ choose a lift in $V$, and an open metric ball
about it whose closure lies in $V$. Images of these balls cover
$D$. A finite subcover and a closed shrinking give compact
subsets $D_i$ covering $D$, each contained in the image of its
chosen ball $U_i$. A closed shrinking is obtained directly by
taking sufficiently small metric neighborhoods in the finite
cover of the compact space and then their closures in $D$.
Make each ball have diameter at most $1/2$, and each $D_i$
have diameter at most $1/2$, refining the finite cover when
needed.

Repeat inside every such pair $(U_i,D_i)$. Lift every point
of $D_i$ inside $U_i$; choose child balls with closures in
$U_i$ and diameter at most $2^{-2}$; use their open images
and a finite closed shrinking to cover $D_i$ by compact child
sets contained in those images and of diameter at most
$2^{-2}$. Continue with bound $2^{-n}$ at level $n$.
Discard empty child sets. This is a finitely branching tree.
Along each infinite branch the closed source balls are nested
and have diameters tending to zero. Completeness gives one
source limit, lying in the initial $V$. The nested nonempty
compact target sets similarly give one target limit in $D$.
At level $n$ that target limit has a lift in the corresponding
source ball, because the whole target set lies in its image.
These lifts converge to the source limit. Continuity of $p$
identifies its image with the target limit.

The infinite-branch space is compact, and the source-limit
map from it is continuous by the source diameter bounds.
Its image $S$ is therefore compact and maps into $D$.
Every $y\in D$ belongs to a nested sequence of child sets,
so some branch limit maps to $y$. This proves $p(S)=D$,
with no continuous endpoint section assumed.

\emph{One compact family, over all selected orbits.}
Fix a complete compatible metric on $L$, and use its uniform
path metric on $\mathscr P(L)$. Choose a compact exhaustion
$E_n$ of the open set $W$, whose interiors cover $W$.
Let $V_n$ be the open neighborhood of the constant identity
path consisting of paths uniformly within $2^{-n}$ of it.
For each $x\in E_n$, openness of (C.64.4) supplies a small
closed coordinate disk $D_{n,x}$ in $e_x(V_n)$ with $x$ in
its interior. These interiors cover $E_n$; retain finitely
many, with centers denoted $x_{n,j}$. The preceding compact
lifting construction gives compact sets
$S_{n,j}\subset V_n$ with
\[
 \{\sigma(1)x_{n,j}:\sigma\in S_{n,j}\}=D_{n,j}.          \tag{C.64.5}
\]
Take all labelled values $\sigma(s)$, for $0\leq s\leq1$
and $\sigma$ in these compact sets, and add the identity.
Their union $\mathcal K$ is compact in $L$: for each $n$
there are finitely many compact value families, and all
families with $n\to\infty$ converge uniformly to the identity.
The same observation proves sequential compactness directly
in the chosen metric. It is path connected, since every
one of its values is connected to the identity by the
corresponding initial segment of its path, entirely in
$\mathcal K$. Compactness and these paths hold for the
actual labelled tuples, including all source lifts and
inverse maps.

For $x$ in the interior of $D_{n,j}$ choose an endpoint
$k_x\in\mathcal K$ with $x=k_xx_{n,j}$. Every other point
of that disk is $k_yx_{n,j}$ for some $k_y\in\mathcal K$.
Consequently
\[
       D_{n,j}\subset\mathcal K\mathcal K^{-1}(x).
                                                               \tag{C.64.6}
\]
Their interiors cover $W$, proving (C.64.1). The projection
of $\mathcal K$ to global plane homeomorphisms is compact
and retains global commutation with $H$ and the fixed set.
We use the same notation for that projected family in the
pointwise estimates below.

\emph{Disks and the exact operator ambiguity.}
For an open orbit with zero-dimensional full local centralizer,
C.30 and C.33 make that centralizer a discrete Lie group
acting freely and properly. The effective image
\[
 \overline H_O=\{h|_O:h\in H_O\}
\]
is a subgroup of this discrete group, and hence is closed
and acts freely and properly. Its actual kernel is exactly
$\{h\in H_O:h|_O=\id\}$, as in (C.30.4).
For any $x\in O$, properness and freeness give a neighborhood
whose closure is disjoint from every nonidentity effective
translate. To see the shrinking explicitly, start with a
compact neighborhood. Properness permits only finitely
many translates meeting it. None of the nonidentity elements
in this finite set fixes $x$; disjoint neighborhoods of their
distinct images allow a further shrinking with disjoint
closures. Intersect also with the interior neighborhood in
(C.64.1). A sufficiently small closed coordinate disk $D$
about $x$ then satisfies (C.64.2) and the disjointness assertion.

Use any family specified in the claim. Distinct representatives
are taken modulo their whole $H$-orbits. An operator carrying
one representative disk to itself stabilizes its source orbit;
the preceding disjointness forces its effective restriction
to be the identity. Thus it fixes that disk pointwise.
For $\psi=(\psi_i)_i$ define
\[
 \mathfrak L(\psi)|_{hD_i}=h\psi_i h^{-1},\qquad
 \mathfrak L(\psi)(x)=x
 \quad\text{off the union of the disk interiors}.        \tag{C.64.7}
\]
Changing $h$ without changing its disk changes it by an
operator acting identically on the representative orbit.
Formula (C.64.7) is therefore unambiguous. All disk boundaries
are fixed; different disk interiors are disjoint. The formula
is a bijection with the identical formula for $\psi_i^{-1}$
as inverse, and it commutes with every element of $H$.

\emph{Uniform control at every fixed frontier.}
Put
\[
          Q=\mathcal K\mathcal K^{-1},\qquad
          P=Q^2=\mathcal K\mathcal K^{-1}
                         \mathcal K\mathcal K^{-1}.      \tag{C.64.8}
\]
These are compact global families, $Q=Q^{-1}$, and every
member fixes the $G$-fixed set and commutes with $H$.
For $x,y\in D_i$ choose $q_x,q_y\in Q$ with
$x=q_xx_i$ and $y=q_yx_i$. Then
$y=q_yq_x^{-1}x\in P(x)$. Commutation with $H$ gives
the same containment for every translated disk. Hence,
including the fixed exterior points by $1\in P$,
\[
 \mathfrak L(\psi)(x)\in P(x),\qquad
 \mathfrak L(\psi)^{-1}(x)\in P(x)                        \tag{C.64.9}
\]
for all disk maps $\psi$ and all $x\in\R^2$.

Let $x_m\to z\in\partial W$, allowing the disk index,
operator translate and every disk map to vary arbitrarily.
Compactness of $P$ and joint continuity of evaluation give
\[
             \sup_{p\in P}|p(x_m)-z|\longrightarrow0,    \tag{C.64.10}
\]
because $p(z)=z$ for every $p\in P$. Equations
(C.64.9)--(C.64.10) prove joint continuity of the pasted
maps and inverses at such points, uniformly over the
entire product. At an interior point of $W$, the stipulated
local finiteness leaves finitely many disk maps in a
neighborhood. Their ordinary disk evaluation and inverse
evaluation are jointly continuous, and they agree with
the identity on their boundaries. The same assertions hold
at exterior interior points, where all maps are the identity.
Thus (C.64.7) gives actual global homeomorphisms and joint
continuity of both families.

Since the plane is locally compact, joint continuity on a
compact set, followed by a finite product-neighborhood
subcover, gives continuity into its compact-open map space.
Apply this also to inverses. This proves continuity of
(C.64.3) for the product topology in its source and the
actual homeomorphism topology in its target. Composition
is componentwise on every disk, so it is a homomorphism.
Restriction to the representative disks is its continuous
inverse on the image, and in particular it is injective.

\emph{The actual contraction.}
Choose once an actual disk chart
$\theta_i:\overline{\mathbb D}\to D_i$. On the closed unit
disk, for $a\in\operatorname{Homeo}(\overline{\mathbb D}
\mathbin{\mathrm{rel}}S^1)$ set
\[
 \mathcal A_s(a)(z)=
 \begin{cases}
 s\,a(z/s),& |z|\leq s,\quad s>0,\\
 z,& |z|\geq s,
 \end{cases}
 \qquad \mathcal A_0(a)=\id.                            \tag{C.64.11}
\]
The two formulas agree on $|z|=s$ because $a$ fixes $S^1$.
Its inverse is $\mathcal A_s(a^{-1})$. At $s=0$ the
displacement of either map is at most $2s$, uniformly in
$a$; at positive $s$ ordinary joint evaluation proves
continuity, also across the joining circle. Therefore
(C.64.11) is a continuous contraction of the disk
homeomorphism group, with $\mathcal A_1(a)=a$ and
$\mathcal A_s(\id)=\id$.
Conjugate this formula by every $\theta_i$, apply it in
each factor with the same $s$, and then apply $\mathfrak L$.
Continuity in the product topology and the proved continuity
of $\mathfrak L$ give the stated contraction of its image.
For a based loop the endpoint identities are preserved by
$\mathcal A_s(\id)=\id$. Formula (C.64.9) controls the
whole contraction and its inverse at every frontier.

This proves actual equivariant disk localization and a
contractible global subgroup, with the original stabilizer
kernel retained. No factorization of an arbitrary original
loop into this subgroup is proved here. Such a controlled
factorization or a simultaneous graph deformation, and
extension at nonfixed frontier points, remain in C.12.
\end{proof}
